For score coordinates \(j=(k,c)\) and \(q=(\ell,d)\), write \(k=(g_k,t_k)\) and \(\ell=(g_\ell,t_\ell)\). By the pairwise-score definition and causal identification, decompose the centered scores as
\[
 \Psi_j=T_j+\sum_{a,s}U_j^{a,s},
\]
where
\[
 T_j=\frac{\mathbf 1\{G_i=g_k,S_i=1\}}{\pi_{g_k,1}}
       \{d_{k,c}(X_i)-\Theta_k(P)\}
\]
and
\[
 U_j^{a,s}=\mathbf 1\{G_i=a,S_i=s\}L_{k,c}^{a,s}(X_i)
       \{\Delta Y_{i,k}-m_{a,s,k}(X_i)\}.
\]
The sums range only over the four cohort-eligibility cells entering the corresponding score. Conditional residual centering gives
\[
 E_P[U_j^{a,s}\mid G_i,S_i,X_i]=0.
\]
Therefore every covariance between a target-cell contrast term and a residual term is zero, including when the residual uses the target cell itself.

For the two target-cell terms, mutual exclusivity of different cohort cells gives \(T_jT_q=0\) if \(g_k\ne g_\ell\). If \(g_k=g_\ell\), iterated expectation gives
\[
 E_P[T_jT_q]
 =E_P\!\left[
 \frac{p_{g_k,1}(X_i)
 \{d_{k,c}(X_i)-\Theta_k(P)\}
 \{d_{\ell,d}(X_i)-\Theta_\ell(P)\}}
 {\pi_{g_k,1}\pi_{g_\ell,1}}
 \right].
\]
This is the first term in the primitive shared-cell formula.

For residual terms, the product \(U_j^{a,s}U_q^{b,r}\) is identically zero unless \((a,s)=(b,r)\), because distinct observed cohort-eligibility cells are mutually exclusive. When the same cell enters both scores, conditioning on that cell and \(X_i=x\) gives
\[
 E_P[U_j^{a,s}U_q^{a,s}\mid X_i=x]
 =p_{a,s}(x)L_{k,c}^{a,s}(x)L_{\ell,d}^{a,s}(x)
   v_{a,s;k\ell}(x).
\]
Summing over precisely the cells shared by the two score coordinates and adding the target-cell term yields exactly \(\Omega^{\mathrm{DDD}}_{jq}\). Since this holds for every pair \((j,q)\),
\[
 \Omega=E_P[\Psi\Psi^{\prime}]=\Omega^{\mathrm{DDD}}.
\]
The same calculation proves the stated interpretation: distinct observed cells contribute no covariance, while two outcome changes observed in the same cell contribute their conditional covariance \(v_{a,s;k\ell}(X_i)\).